\section{软件实现思路}

\subsection{逆运动学：从坐标到关节角}

运动控制接口支持末端坐标和关节角两种输入。对于末端坐标 $(x,y,z)$，逆运动学模块需要计算底座、肩、肘和腕关节的目标角度；夹爪开合量由独立参数给出。该结构可分解为底座水平旋转和平面连杆求解两部分：

\begin{enumerate}
  \item 底座角由目标点在水平面的投影确定：
  \[
    \theta_0=\operatorname{atan2}(y,x).
  \]
  \item 在底座旋转后的竖直平面内，先固定或给定末端俯仰角，再由腕部目标位置求解肩、肘两连杆的关节角，最后补偿腕关节角。
\end{enumerate}

设大臂、前臂和腕部等效长度分别为 $l_1,l_2,l_3$，末端目标俯仰角为 $\phi$，底座轴到肩关节的高度为 $z_0$。令
\[
  r=\sqrt{x^2+y^2},\qquad
  r_w=r-l_3\cos\phi,\qquad
  z_w=z-z_0-l_3\sin\phi ,
\]
则肘关节可由余弦定理求得：
\[
  c_2=\frac{r_w^2+z_w^2-l_1^2-l_2^2}{2l_1l_2},\qquad
  \theta_2=\operatorname{atan2}\!\left(\pm\sqrt{1-c_2^2},c_2\right).
\]
肩关节和腕关节再由
\[
  \theta_1=\operatorname{atan2}(z_w,r_w)
  -\operatorname{atan2}(l_2\sin\theta_2,l_1+l_2\cos\theta_2),
  \qquad
  \theta_3=\phi-\theta_1-\theta_2
\]
计算。程序同时检查 $|c_2|\leq1$、关节限位和肘上/肘下两组解，优先选择与当前姿态变化较小且不碰撞的一组。靠近全伸直奇异位形时限制目标半径，避免微小坐标噪声引起关节角突变。

几何解计算量小，也便于在求解前检查工作空间和关节限位，因此首版不采用雅可比迭代。对于不可达目标，控制器不会直接执行未经约束的角度，而是将目标投影到允许工作空间的边界，并向上位机返回钳制标志。各关节还需分别标定中位脉宽、方向、有效角度范围和机械零点。为减小舵机同时启动造成的冲击，关节目标之间采用限速插值；需要越过障碍或抓取物体时，使用“抬升到安全高度、平移、下降”的离散路径，而不假设关节空间直线始终无碰撞。

\subsection{视觉管线：一条流水线五种模式}

摄像头图像由 DMA 写入缓冲区，视觉任务取得完整帧后，根据当前模式调用对应的检测模块。各模式共用图像采集、结果滤波和事件上报代码，检测模块只输出目标位置、有效标志和时间戳。考虑到 esp32-camera 官方文档指出 YUV/RGB 帧会增加 PSRAM 带宽压力，首版在 QQVGA 下使用单帧缓冲，并针对模式选择输入格式，避免同时保留多份全尺寸彩色图像。工作缓冲区在启动时分配并循环使用，运行过程中不频繁申请内存。

当前的算法选型如下：

\begin{itemize}
  \item \textbf{人脸检测}：基于 ESP-DL/ESP-WHO 的轻量模型，先在自配置控制板上完成摄像头驱动和板级适配，再根据端到端采集及推理耗时确定输入尺寸与检测周期；
  \item \textbf{指尖检测}：YCbCr 肤色二值化、形态学开运算、最大连通域筛选和凸包顶点分析；
  \item \textbf{颜色物体检测}：HSV 阈值分割、形态学去噪和最大连通域质心计算；
  \item \textbf{手势识别}：在肤色轮廓上计算 solidity、凸包缺陷数和长宽比，并通过连续 3 帧投票确认类别。
\end{itemize}

对人脸追随而言，控制器可以直接使用目标相对图像中心的偏差；对桌面指尖和物体跟踪，则需要建立图像平面到桌面平面的映射。摄像头必须固定在不随关节运动的底座支架上，再在桌面选取分布充分的已知坐标点，利用单应性矩阵完成平面标定。该方法只对已标定平面有效，目标高度变化不能由单目图像直接恢复，需要使用预设高度或额外传感信息。Python SDK 将提供标定数据采集、重投影误差检查和参数写入功能；摄像头支架一旦移动，必须重新标定。

\subsection{JSON 协议：固件和上位机的共同语言}

串口协议运行在 ESP32-S3 原生 USB 串口链路上，采用一行一条 JSON 消息的格式。每个请求包含协议版本、消息类型、序号、参数和可选超时时间；固件响应中返回相同序号、执行状态、结果或错误码。目标出现、目标丢失和动作完成等异步事件使用独立的事件类型，由固件主动发送。

首版指令集覆盖末端移动、关节角控制、夹爪开合、视觉模式切换、单次检测、动作播放、状态查询和紧急停止。协议解析器限制单行长度，并检查字段类型和参数范围，避免畸形消息占用过多内存或绕过关节限位。所有运行期日志也封装为 JSON 事件，并通过同一个发送队列串行输出，防止多任务同时写串口造成半行消息交错。ESP32-S3 复位时原生 USB 串口可能输出 ROM 启动信息，SDK 因此在握手成功前忽略非 JSON 行。文本格式便于调试，性能不足时再评估二进制协议。

\subsection{双核任务划分}

ESP-IDF 的 FreeRTOS 支持为任务设置核心亲和性。系统将周期控制任务与计算量较大的视觉任务固定到不同 CPU 核，但不把核编号写死在算法接口中，最终分配以运行时间统计和相机驱动测试为准：

\begin{itemize}
  \item \textbf{控制任务}：以 50\,Hz 周期更新舵机目标，完成逆运动学、轨迹插值、关节限位和急停处理；
  \item \textbf{感知与通信任务}：运行视觉检测、OLED 刷新和 JSON 协议收发，视觉处理目标为 12--15\,FPS。
\end{itemize}

两个核之间通过 FreeRTOS 队列传递带时间戳的目标和事件，目标队列只保留最新值，避免视觉处理短时拥塞后机械臂追赶过期位置。PCA9685 和 MPU6050 共享 I²C 时，由单一总线服务任务串行执行事务，其他任务通过队列提交请求，避免跨核并发访问驱动。该划分可以降低感知任务阻塞对舵机更新周期的影响，但仍需通过看门狗、优先级设置和运行时间统计验证周期，而不能仅凭“双核”推断实时性。

\subsection{动作库}

动作库计划包含挥手、点头、摇头等预置动作。每套动作由若干关节关键帧和持续时间组成，执行前统一进行角度检查，帧间采用关节空间线性插值并限制最大角速度。动作文件在 PC 端以 JSON 保存，经 SDK 校验后发送给固件，避免把大量可变动作固化在程序中。

\subsection{安全措施}

桌面机械臂速度和负载较低，但仍需防范夹伤、堵转、结构碰撞、掉电下坠和供电异常。首版安全设计包括：

\begin{itemize}
  \item \textbf{指令约束}：所有运动指令在执行前检查关节角、末端工作空间、最大角速度和单次动作时长；
  \item \textbf{自碰撞限制}：根据结构尺寸建立相邻关节的组合限位，禁止前臂与大臂、夹爪与底座进入明显碰撞区域；
  \item \textbf{堵转防护}：MG90S 不提供外部位置反馈，无法仅凭角度指令判断堵转。因此首版采用限速、动作超时和夹爪行程限制；若实测堵转问题明显，再增加分路或总线电流检测；
  \item \textbf{输出使能}：PCA9685 的 \texttt{OE} 由硬件上拉保持默认关闭，固件完成零位和限位加载后才允许输出。软件异常时可快速撤销使能；
  \item \textbf{紧急停止}：软件急停先停止新轨迹并撤销 PWM 输出，另设置可直接断开舵机 5\,V 电源的实体开关。舵机失电后关节可能因重力下落，因此肩关节采用机械限位或弹性配重，并限制人员靠近夹爪工作区。恢复前需重新标定并进入待机状态。
\end{itemize}
